\section{Introduction}\label{sec:intro}

Autoregressive language models sometimes fall into repetition loops: the model emits a short span and then reproduces it until the sample is dominated by a fixed cycle. Most remedies are decoding rules applied one token at a time. They look at the next-token distribution and at a short history, and they penalize or forbid candidates that would repeat recent text. A common design fixes how much repetition risk a single step may carry, a \emph{per-step risk budget}, and re-solves it at every step. Whether a loop occurs, however, is a property of the whole sequence. The distortion that a rule introduces, measured here as the KL divergence of the controlled generation from the model's own, is spent at every step whether or not a loop is at stake.

We ask one question: \emph{why can a fixed per-step risk budget waste KL divergence when the goal is to prevent a loop anywhere in the generated sequence, and by how much?} The answer we give is not that token-by-token control is inefficient in itself. A one-step controller has two dual parameters, a risk budget and its Lagrange multiplier, the \emph{price} of risk. Holding the budget fixed and holding the price fixed are different designs. We show that the cheapest way to lower the probability of a loop holds the price constant across histories, and that a fixed budget instead induces prices that vary by orders of magnitude, which is where the KL is wasted.

\subsection{Contributions}
\begin{enumerate}[leftmargin=*,itemsep=2pt]
\item \textbf{A mechanism with an explicit waste law (Section~\ref{sec:theory}).} The probability of a loop falls by the expected sum of per-step hazard reductions, and the minimum KL for a given reduction is known in closed form. To second order, a controller that reduces the hazard by $\Delta(h)$ at history $h$ costs a factor $1+\mathrm{CV}^2$ more than the optimum, where $\mathrm{CV}$ is the coefficient of variation of its implied per-step prices; a fixed budget costs at least the reciprocal of the share of the event's variance on the histories it touches. These are elementary consequences of known tools, which we credit in Section~\ref{sec:related}; what we add is their use as an accounting of repetition control.
\item \textbf{An exact test, with infeasibility separated from allocation (Section~\ref{sec:small}).} In finite-state systems where the optimum, the controllers and the loop probabilities are computed exactly, a fixed price stays within 13\% of the minimum KL in every setting, whereas a fixed budget costs between 1.3 and more than 200 times the minimum. The excess is present with no infeasible budget at all, and removing infeasibility does not remove it.
\item \textbf{A qualified test on a language model (Section~\ref{sec:real}).} On a character-level model with a validated loop detector we check the prediction that a fixed-price rule reaches a given loop rate at lower KL than a fixed-budget rule, and we state which findings of the exact system the language-model data support and which remain hypotheses.
\end{enumerate}

\subsection{Scope}
Recurrence-risk decoding is used here as an instrument for studying the question, not proposed as a practical decoder. Observed loop counts are counts at one temperature on a few models, not guarantees. KL divergence is a measure of distortion, not of readability or text quality. The language-model estimates are descriptive; they do not say how much distortion a practical decoder could save. Models beyond the primary one, the loop-detector validation, baseline details, a pilot human evaluation, a rollout-based hazard estimate and the implementation defects found along the way are in the appendices.
